\section{Comparison certificates and bounded-loss calibration}
\label{app:witness-analysis}

This appendix specifies the stronger comparison bounds used in the additional
experiments. Each formula concerns independent, identically distributed
calibration units and a model
class fixed before calibration. The concentration arguments are established
tools; the comparison evaluates their informativeness and computational cost
against the pair witness.

\subsection{Relation to an inverse-contrastive population bound}
\label{app:inverse-contrastive}

Consider the same scalar target law and an independent, identically
distributed pair, with $\mathbb E Y^2<\infty$ and $\kappa>0$.
For a threshold $a>0$, define
$F=\{|Y-Y'|>a\}$, $\rho=\Pr(F)$, and
$M=\mathbb E[\mathbf1_F d(U,U')^2]\leq\varepsilon<\infty$.
Put $v=(|Y-Y'|-\kappa d(U,U'))_+$.
Since $\mathbf1_F|Y-Y'|\leq\mathbf1_Fv+
\kappa\mathbf1_Fd(U,U')$, the $L^2$ triangle inequality gives
\begin{equation}
 \begin{split}
 W_\kappa&\geq\tfrac12\|\mathbf1_Fv\|_{L^2}^2
 \geq\tfrac12(a\sqrt\rho-\kappa\sqrt M)_+^2\\
 &\geq\tfrac12(a\sqrt\rho-\kappa\sqrt\varepsilon)_+^2.
 \end{split}
 \label{eq:inverse-contrastive-comparison}
\end{equation}
The inverse-contrastive-loss hypothesis of
\citet[Lemma~2 and Appendix~7.3]{akash2021inverse} bounds $M$ by
$\varepsilon$. In their continuous-target setting and positive regime
$\varepsilon<a^2\rho^2/\kappa^2$, their displayed population lower bound
is $(a\rho-\kappa\sqrt\varepsilon)^2/4$; it is no larger than
\eqref{eq:inverse-contrastive-comparison} because $\sqrt\rho\geq\rho$.
The paired-residual principle therefore has a direct antecedent. The
positive-part functional and conditional independence sharpen its population
consequence. This comparison does not assert dominance of finite-sample
confidence bounds, whose corrections and sampling laws must also agree.

\subsection{A range-aware centered lower bound}
\label{app:centered-range-aware}

For this scalar comparator, let $U,Y\in\mathbb R$ have known interval
ranges of lengths $D_U,D_Y$, and use $d(u,u')=|u-u'|$.
For $n\geq2$ independent, identically distributed pairs $(U_i,Y_i)$, write
$s_U,s_Y$ for the square roots of the unbiased sample variances and
$\widehat v_Y=n^{-1}\sum_i(Y_i-\bar Y)^2$ for the target variance with
denominator $n$. Define
$c_\alpha=\sqrt{2\log(1/\alpha)/(n-1)}$.
For a total failure probability $\eta\in(0,1)$, use
$\alpha=\eta/3$ if $D_U>0$ and $\alpha=\eta/2$ if $D_U=0$, and set
\begin{align}
 \sigma_Y^-&=\max\left\{(s_Y-D_Yc_\alpha)_+,\,
 \sqrt{\left(\widehat v_Y-D_Y^2
                    \sqrt{\frac{\log(1/\alpha)}{2n}}\right)_+}\right\},
 \nonumber\\
 \sigma_U^+&=
 \begin{cases}
   \min\{s_U+D_Uc_\alpha,D_U/2\},&D_U>0,\\
   0,&D_U=0,
 \end{cases}
 \qquad
 L_{\rm ctr}=(\sigma_Y^--\kappa\sigma_U^+)_+^2.
 \label{eq:range-aware-centered}
\end{align}
Then $L_{\rm ctr}\leq R^*_\kappa$ with probability at least $1-\eta$.
The maximum uses separately allocated confidence statements; it is not an
unallocated choice of the more favorable bound.

To verify this, rescale \citet[Theorem~10]{maurer2009empirical} to obtain
the one-sided events
$\sqrt{\operatorname{Var}(Y)}\geq s_Y-D_Yc_\alpha$ and
$\sqrt{\operatorname{Var}(U)}\leq s_U+D_Uc_\alpha$.
Also, $\widehat v_Y=\inf_a n^{-1}\sum_i(Y_i-a)^2$, where $a$ may be
restricted to the known target interval. Replacing one observation changes
this infimum by at most $D_Y^2/n$, and
$\mathbb E\widehat v_Y=(1-1/n)\operatorname{Var}(Y)
\leq\operatorname{Var}(Y)$.
The bounded-differences inequality supplies the second target lower bound.
The deterministic interval bound
$\operatorname{Var}(U)\leq D_U^2/4$ justifies the cap.
Union allocation and \eqref{eq:pair-dominance} with constant $A$ finish
the argument. When the representation is structurally constant, its variance
is exactly zero and requires no confidence correction. A zero sample variance
alone does not justify that substitution.

\subsection{One-sided relative-entropy bounds}
\label{app:kl-calibration}

For $q,p\in[0,1]$, define binary relative entropy
\[
 \operatorname{kl}(q\Vert p)
 =q\log(q/p)+(1-q)\log\{(1-q)/(1-p)\},
\]
using its continuous endpoint values and $0\log0=0$.
For $c>0$, let $\operatorname{kl}_-(q,c)$ and
$\operatorname{kl}_+(q,c)$ be the smallest and largest $p\in[0,1]$
satisfying $\operatorname{kl}(q\Vert p)\leq c$.
They lie respectively below and above $q$ and can be found by scalar
bisection. In particular,
$\operatorname{kl}_-(0,c)=0$ and
$\operatorname{kl}_+(1,c)=1$.

Let $V_1,\ldots,V_n$ be independent and identically distributed in
$[0,B]$, with $B>0$, mean $\mu$, and sample mean $\bar V$.
For $\eta\in(0,1)$ and $c=\log(1/\eta)/n$, each of
\begin{equation}
 B\operatorname{kl}_-(\bar V/B,c)\leq\mu,
 \qquad
 \mu\leq B\operatorname{kl}_+(\bar V/B,c)
 \label{eq:kl-bounded-mean}
\end{equation}
has failure probability at most $\eta$. A joint two-sided claim allocates
failure probability to both directions. If $B=0$, both endpoints are zero.
The scores need not be Bernoulli observations.

For completeness, convexity of the exponential gives, for every real $t$,
\[
 \mathbb E e^{tV_1/B}
 \leq 1-\mu/B+(\mu/B)e^t.
\]
Independence, the exponential Markov inequality, and minimization over
$t\geq0$ or $t\leq0$ give the bounded-variable Chernoff bounds
\citep{hoeffding1963}. Their monotone inversion yields
\eqref{eq:kl-bounded-mean}.
Applying the upper endpoint to losses of a frozen forecaster gives an
alternative to \eqref{eq:risk-ucb}, under the same sampling and range
assumptions. An episode average is one bounded score. Choosing between
different concentration methods after seeing calibration requires a joint
allocation, just as in Section~\ref{sec:theory}.

\subsection{A certified empirical-minimum comparator}
\label{app:erm-comparator}

An empirical-minimum lower bound provides a comparison with structural
certificates that avoid optimization on calibration observations
\citep[Section~2.5]{kumar2026limits}. Let $\mathcal G$ be a fixed nonempty
class with losses $\ell_g(X,Y)\in[0,B]$, $B>0$, and define
the following quantities from independent, identically distributed
observations $(X_i,Y_i)$ of the target law:
\[
 R^*=\inf_{g\in\mathcal G}\mathbb E\ell_g(X,Y),\qquad
 a_n=\inf_{g\in\mathcal G}\frac1n\sum_i\ell_g(X_i,Y_i).
\]
Assume these infima define measurable statistics. If a computed value
$d_n$ satisfies $d_n\leq a_n$, put $q_n=\max\{0,d_n/B\}$.
Since $a_n\leq B$, one has $q_n\in[0,1]$. Then
\begin{equation}
 L_{\rm ERM}
 =B\operatorname{kl}_-\!\left(q_n,\frac{\log(1/\eta)}{n}\right)
 \leq R^*
 \label{eq:erm-kl-lower}
\end{equation}
with probability at least $1-\eta$.

Indeed, $a_n\leq n^{-1}\sum_i\ell_g(X_i,Y_i)$ for every fixed $g$.
For $t\geq0$, the preceding moment-generating-function bound therefore gives
\[
 \mathbb E e^{tn a_n/B}
 \leq\inf_{g\in\mathcal G}
       \{1-R(g)/B+(R(g)/B)e^t\}^{n}
 =\{1-R^*/B+(R^*/B)e^t\}^{n},
\]
where $R(g)=\mathbb E\ell_g(X,Y)$.
The same upper-tail inversion proves \eqref{eq:erm-kl-lower}, without
requiring an attained population optimum. This direction needs no uniform
complexity penalty: the empirical infimum lies below the empirical risk
of every fixed population competitor. The result lower-bounds the risk of
every class member; it supplies neither an upper risk guarantee nor an
excess-risk bound for a predictor fitted to these calibration observations.

For scalar Lipschitz regression on a finite calibration sample, the empirical
infimum is a convex quadratic program
\citep[Section~III-A]{gottlieb2017regression}. Restricting predictions to the known
target interval preserves the optimum, because clipping reduces loss and
preserves the Lipschitz budget. Let $u_1,\ldots,u_m$ be its distinct
observed representation values and $v_j$ their fitted outputs in the known
interval $[y_-,y_+]$. Set
\[
 w_j=\frac{\#\{i:U_i=u_j\}}n,\qquad
 t_j=\frac1n\sum_{i:U_i=u_j}Y_i,\qquad
 \nu=\frac1n\sum_iY_i^2.
\]
Here $\#$ counts sample indices. Minimize
$\nu+\sum_j(w_jv_j^2-2t_jv_j)$ subject to
$v_j-v_k\leq\kappa d(u_j,u_k)$ for all ordered pairs and
$y_-\leq v_j\leq y_+$. The scalar extension used in
Proposition~\ref{prop:stratified-exactness}, followed by interval projection,
shows that these finite constraints describe exactly the empirical values
of the global Lipschitz class. For ordered scalar inputs with absolute
distance, the adjacent constraints imply all pair constraints by summation.

Write the $L$ pairwise and interval constraints as $Av\leq b$, with
$A\in\mathbb R^{L\times m}$ and $b\in\mathbb R^L$.
For any vector $\lambda\in[0,\infty)^L$ of constraint multipliers,
a valid dual lower value is
\begin{equation}
 d(\lambda)=\nu-\lambda^\top b-
 \sum_{j=1}^m\frac{((A^\top\lambda)_j-2t_j)^2}{4w_j}.
 \label{eq:empirical-risk-dual}
\end{equation}
This minimizes the Lagrangian over all $v\in\mathbb R^m$; its minimizing
coordinate is $(2t_j-(A^\top\lambda)_j)/(2w_j)$.
The observed-node convention ensures $w_j>0$.
Weak duality gives $d(\lambda)\leq a_n$.
A primal objective alone is an upper bound on $a_n$ and cannot replace this
dual value. The implementation records primal feasibility and the primal--dual
gap; its floating-point checks do not constitute interval-arithmetic proofs.
